\chapter{引言}
\label{cha:intro}

引言（或绪论）简要说明研究工作的目的、范围、相关领域的前人工作和知识空白、理论基础和分析、研究设想、研究方法和实验设计、预期结果和意义等。应言简意赅，不要与摘要雷同，不要成为摘要的注释。一般教科书中有的知识，在引言中不必赘述。

\section{背景}

学位论文为了需要反映出作者确已掌握了坚实的基础理论和系统的专门知识，具有开阔的科学视野，对研究方案作了充分论证，因此，有关历史回顾和前人工作的综合评述，以及理论分析等，可以单独成章，用足够的文字叙述。正文是学位论文的核心部分，占主要篇幅，可以包括：调查对象、实验和观测方法、仪器设备、材料原料、实验和观测结果、计算方法和编程原理、数据资料、经过加工整理的图表、形成的论点和导出的结论等。


由于研究工作涉及的学科、选题、研究方法、工作进程、结果表达方式等有很大的差异，对正文内容不能作统一的规定。但是，必须实事求是，客观真切，准确完备，合乎逻辑，层次分明，简练可读。

由于研究工作涉及的学科、选题、研究方法、工作进程、结果表达方式等有很大的差异，对正文内容不能作统一的规定。但是，必须实事求是，客观真切，准确完备，合乎逻辑，层次分明，简练可读。具体写作规范参考《云南大学本科毕业论文（设计）写作规范》\footnote{《云南大学本科毕业论文（设计）写作规范》：\url{https://co2.cnki.net/Notice.html?typeId=3596\&dp=ynu\&r=1776851073592}}，以及《云南大学研究生学位论文写作规范》\footnote{《云南大学研究生学位论文写作规范》：\url{http://www.grs.ynu.edu.cn/info/1036/1352.htm}}。




\section{二级标题}

命令：\textbackslash section\{二级标题\}。由于研究工作涉及的学科、选题、研究方法、工作进程、结果表达方式等有很大的差异，对正文内容不能作统一的规定。但是，必须实事求是，客观真切，准确完备，合乎逻辑，层次分明，简练可读。



\subsection{三级标题}

命令：\textbackslash section\{三级标题\}。实事求是，客观真切，准确完备，合乎逻辑，层次分明，简练可读。

\subsubsection{四级标题}

命令：\textbackslash subsubsection\{四级标题\}。必须实事求是，客观真切，准确完备，合乎逻辑，层次分明，简练可读。
